% Generated by Listas/make_statements.py from pratica-lecture-07.md, item 6. Do not edit.
\textbf{Source.} Muñoz-Garcia \& Toro-Gonzalez (2019), \emph{Strategy and Game Theory: Practice Exercises with Answers}, Ch. 8, Exercise 4 ``Second Price Auction'', p.~362 {[}PDF p.~372{]}. Reproduced verbatim (statement only; the book's worked answer follows it). \textbf{Answer in the source:} yes, the book's worked answer follows the exercise.



\begin{quote}

Consider a setting where every bidder privately observes his/her own valuation for the object. In addition, assume that \(N\) bidders interact in an auction in which they simultaneously and independently submit their own bids, the highest bidder wins the object, but he/she pays the second-highest bid (this is the so-called second-price auction, SPA). All other bidders pay zero.



\begin{enumerate}

\def\labelenumi{\alph{enumi}.}

\item

  Argue that it is a weakly dominant strategy for a bidder to bid her value in the SPA.

\item

  Are optimal bids affected if the joint distribution of bidders' values exhibits correlation, i.e., if bidder \(i\) observes that he has a high valuation for the object he can infer that other bidders must also have a high valuation (in the case of positively correlated values) or that, instead, they must have a low valuation for the object (in the case of negatively correlated values)?

\item

  Are optimal bids affected by the number of bidders participating in the auction? Are optimal bids affected by bidders' risk aversion?

\end{enumerate}

\end{quote}

